\documentclass[10pt,a4paper]{article}

\usepackage[margin=0.62in]{geometry}
\usepackage[T1]{fontenc}
\usepackage[utf8]{inputenc}
\usepackage{lmodern}
\usepackage{microtype}
\usepackage[dvipsnames]{xcolor}
\usepackage{enumitem}
\usepackage{tabularx}
\usepackage{titlesec}
\usepackage[hidelinks]{hyperref}
\usepackage{fancyhdr}

\definecolor{accent}{HTML}{17365D}
\hypersetup{
  colorlinks=true,
  urlcolor=accent,
  linkcolor=accent,
  pdfauthor={Xijia Tao},
  pdftitle={Xijia Tao -- Curriculum Vitae}
}

\pagestyle{fancy}
\fancyhf{}
\renewcommand{\headrulewidth}{0pt}
\cfoot{\small\color{gray}\thepage}

\setlength{\parindent}{0pt}
\setlength{\parskip}{0pt}
\setlist[itemize]{leftmargin=1.25em,itemsep=1pt,topsep=2pt,parsep=0pt}
\titleformat{\section}
  {\large\bfseries\color{accent}}
  {}{0pt}{}[\vspace{-0.35em}\color{accent}\titlerule]
\titlespacing*{\section}{0pt}{0.75em}{0.4em}

\newcommand{\cvheading}[4]{%
  \textbf{#1}\hfill #2\\[-1pt]
  \textit{#3}\hfill\textit{#4}\par
}
\newcommand{\me}{\textbf{Xijia Tao}}
\newcommand{\paper}[4]{%
  \textbf{#1} \hfill #2\\[-1pt]
  #3\\[-1pt]
  \textit{#4}\par\vspace{0.15em}
}

\begin{document}

\begin{center}
  {\LARGE\bfseries Xijia Tao}\\[3pt]
  Hong Kong SAR \enspace$\vert$\enspace
  \href{mailto:xjtao@cs.hku.hk}{xjtao@cs.hku.hk} \enspace$\vert$\enspace
  \href{mailto:ciellltao@gmail.com}{ciellltao@gmail.com}\\[2pt]
  +852 5467 9633 \enspace$\vert$\enspace +86 186 4115 5390 \enspace$\vert$\enspace
  \href{https://xijia-tao.github.io}{Webpage} \enspace$\vert$\enspace
  \href{https://scholar.google.com/citations?hl=en\&user=PeIV968AAAAJ}{Google Scholar} \enspace$\vert$\enspace
  \href{https://github.com/xijia-tao}{GitHub} \enspace$\vert$\enspace
  \href{https://www.linkedin.com/in/xijia-ciel-tao}{LinkedIn} \enspace$\vert$\enspace
  \href{https://x.com/xijia_tao}{X}
\end{center}

\section{Research Interests}

Multimodal deep research and browsing agents; diffusion language models; large language model reasoning and safety.

\section{Education}

\cvheading{The University of Hong Kong}{Hong Kong SAR}
  {Ph.D. in Computer Science, HKU NLP Lab}{Sep 2024 -- Jun 2028 (expected)}

\vspace{0.35em}
\cvheading{The University of Hong Kong}{Hong Kong SAR}
  {B.Eng. in Computer Science}{Sep 2020 -- Jun 2024}
\begin{itemize}
  \item Full-tuition scholarship for four years of undergraduate study; Dean's Honours List, Faculty of Engineering (2021--22).
\end{itemize}

\section{Research Experience}

\cvheading{HKU NLP Lab, The University of Hong Kong}{Hong Kong SAR}
  {Ph.D. Student}{Sep 2024 -- Present}
\begin{itemize}
  \item Develop benchmarks and methods for provenance-aware multimodal browsing agents, including MMSearch-Plus (ICLR 2026).
  \item Study training and inference for diffusion language models, including adaptive decoding and trajectory-aware Boltzmann modeling.
  \item Investigate behavioral compression of agent skills for contextual adaptation.
\end{itemize}

\vspace{0.25em}
\cvheading{Zhipu AI}{}
  {AI Research Intern}{Jul 2024 -- Oct 2024}
\begin{itemize}
  \item Developed methods for synthesizing domain-specific training data and evaluated data quality on reasoning benchmarks by training GLM-4-9B models.
\end{itemize}

\vspace{0.25em}
\cvheading{HKU NLP Lab, The University of Hong Kong}{Hong Kong SAR}
  {Research Assistant}{Sep 2021 -- Mar 2024}
\begin{itemize}
  \item Researched vision-language model safety, multilingual transfer, knowledge distillation, and pretrained language models for the legal domain.
\end{itemize}

\vspace{0.25em}
\cvheading{Huawei Hong Kong Research Center}{Hong Kong SAR}
  {Machine Learning Research Intern}{Jun 2023 -- Sep 2023}
\begin{itemize}
  \item Implemented multimodal networks in MindSpore for CUDA and Ascend environments and studied language--vision feature alignment for large generative models.
\end{itemize}

\section{Publications}

\paper{SoftSkill: Behavioral Compression for Contextual Adaptation}{2026}
  {\me, Yihua Teng, Xinyu Fu, Ziru Liu, Kecheng Chen, Yuzhi Zhao, Suiyun Zhang, Rui Liu, Lingpeng Kong}
  {arXiv:2606.20333. \href{https://arxiv.org/abs/2606.20333}{Paper}}

\paper{Self-Distilled Trajectory-Aware Boltzmann Modeling: Bridging the Training--Inference Discrepancy in Diffusion Language Models}{2026}
  {Kecheng Chen, Ziru Liu, \me, Hui Liu, Yibing Liu, Xinyu Fu, Shi Wu, Suiyun Zhang, Dandan Tu, Lingpeng Kong, Rui Liu, Haoliang Li}
  {arXiv:2605.11854. \href{https://arxiv.org/abs/2605.11854}{Paper}}

\paper{MMSearch-Plus: Benchmarking Provenance-Aware Search for Multimodal Browsing Agents}{2026}
  {\me, Yihua Teng, Xinxing Su, Xinyu Fu, Jihao Wu, Chaofan Tao, Ziru Liu, Haoli Bai, Rui Liu, Lingpeng Kong}
  {The Fourteenth International Conference on Learning Representations (ICLR 2026). \href{https://mmsearch-plus.github.io/}{Project} $\vert$ \href{https://arxiv.org/abs/2508.21475}{Paper} $\vert$ \href{https://github.com/mmsearch-plus/MMSearch-Plus}{Code}}

\paper{Beyond Confidence: Adaptive and Coherent Decoding for Diffusion Language Models}{2026}
  {Kecheng Chen, Ziru Liu, \me, Hui Liu, Xinyu Fu, Suiyun Zhang, Dandan Tu, Lingpeng Kong, Rui Liu, Haoliang Li}
  {The Forty-third International Conference on Machine Learning (ICML 2026). \href{https://tonyckc.github.io/CCD-DLM-Project/}{Project} $\vert$ \href{https://arxiv.org/abs/2512.02044}{Paper}}

\paper{OmniPlay: Benchmarking Omni-Modal Models on Omni-Modal Game Playing}{2025}
  {Fuqing Bie, Shiyu Huang, \me, Zhiqin Fang, Leyi Pan, Junzhe Chen, Min Ren, Liuyu Xiang, Zhaofeng He}
  {arXiv:2508.04361. \href{https://arxiv.org/abs/2508.04361}{Paper}}

\paper{ImgTrojan: Jailbreaking Vision-Language Models with ONE Image}{2025}
  {\me, Shuai Zhong, Lei Li, Qi Liu, Lingpeng Kong}
  {Proceedings of NAACL 2025. \href{https://aclanthology.org/2025.naacl-long.360/}{Paper} $\vert$ \href{https://github.com/xijia-tao/ImgTrojan}{Code}}

\paper{Language Versatilists vs. Specialists: An Empirical Revisiting on Multilingual Transfer Ability}{2023}
  {Jiacheng Ye, \me, Lingpeng Kong}
  {arXiv:2306.06688. \href{https://arxiv.org/abs/2306.06688}{Paper}}


\end{document}
